\section{Experiments}
\label{sec:experiments}

\subsection{Setup}

\paragraph{Reference SSD evaluation.}
The original S$^2$D paper evaluates batch-size-one greedy decoding with the
target on $4\times$H100 and, for SSD, the asynchronous draft on a separate
H100~\cite{saguaro2024}. Its Llama setup uses Llama-3.1-70B-Instruct with a
Llama-3.2-1B draft model; the Qwen appendix uses Qwen3-32B with a Qwen3-0.6B
draft model~\cite{yang2025qwen3}. The benchmark pool contains $128$ prompts
from each of HumanEval, Alpaca, GSM8K, and UltraFeedback, for $512$ prompts
total, with $512$ output tokens per prompt.

\paragraph{Our local evaluation.}
Our diffusion experiments were constrained by available compute. We used
shared A6000-class machines for implementation and smoke testing, and several
larger runs were blocked by other processes occupying most GPU memory. The
reported prototype numbers should therefore be read as diagnostic
measurements, not as final throughput claims. This limitation matters:
S$^2$D is a latency-hiding method, so evaluating it on a small verifier or on
busy GPUs can make communication and scheduler overhead look much worse than
it would with a larger target.

\paragraph{Models.}
For the traditional block-diffusion path, we used the small Qwen-family stack
available in the repository and the block/MDLM-style draft checkpoints
supported by our backend. For the DFlash hybrid, we used Qwen3-8B as the
target, Qwen3-0.6B as the normal S$^2$D AR drafter, and
\texttt{z-lab/Qwen3-8B-DFlash-b16} as the target-conditioned DFlash drafter.
The DFlash checkpoint has block size $16$, so all DFlash runs use $k=15$.

\paragraph{Metrics.}
We report end-to-end generated-token throughput, average accepted length per
verify step, S$^2$D cache hit rate, DFlash seed latency, and SSD tree-decode
latency. Throughput is the user-visible metric. Acceptance and cache hit rate
diagnose draft quality and S$^2$D prediction quality. Seed and tree-decode
latencies identify whether the bottleneck is DFlash, the AR tree cache, or
cross-process orchestration.

\subsection{Reference numbers from SSD and DFlash}

Table~\ref{tab:ssd-reference} summarizes the published S$^2$D numbers most
relevant to this project. The Qwen3-32B results are especially important: they
show that S$^2$D still improves over ordinary SD, but the speedup over AR is
smaller than in the Llama-70B setting because the verifier is faster and the
draft/communication overheads are less easily amortized.

\begin{table}[t]
  \caption{Published S$^2$D results from Kumar et al.~\cite{saguaro2024}.
  Values are end-to-end decode throughput in tokens/s averaged over four
  datasets.}
  \label{tab:ssd-reference}
  \centering
  \footnotesize
  \begin{tabular}{@{}lrrrr@{}}
    \toprule
    Model pair & AR & SD & SSD & SSD/AR \\
    \midrule
    Llama 70B / 1B & 54.7 & 161.8 & 255.8 & $4.68\times$ \\
    Qwen3 32B / 0.6B & 88.8 & 136.8 & 203.8 & $2.29\times$ \\
    \bottomrule
  \end{tabular}
\end{table}

Table~\ref{tab:dflash-reference} gives the DFlash reference point. DFlash is
not merely a generic block-diffusion model; it uses target activations and
generates the whole block in one conditioned draft pass. Published SGLang
results on a single B200 show large Qwen3-8B speedups, which is why DFlash is
a more promising diffusion substrate for S$^2$D than iterative block
diffusion.

\begin{table}[t]
  \caption{Published DFlash SGLang throughput for Qwen3-8B on a single B200
  with FlashAttention-4~\cite{chen2026dflash}.}
  \label{tab:dflash-reference}
  \centering
  \footnotesize
  \begin{tabular}{@{}lrrrr@{}}
    \toprule
    Task & AR & DFlash & Speedup & Avg. $\tau$ \\
    \midrule
    Math500 & 230 & 1175 & $5.1\times$ & 8.01 \\
    HumanEval & 229 & 955 & $4.2\times$ & 6.50 \\
    \bottomrule
  \end{tabular}
\end{table}

\subsection{Our findings}

\paragraph{Traditional block diffusion is too slow for SSD.}
The direct block-diffusion integration did not produce useful S$^2$D speedups.
The reason is mechanical: although a diffusion forward predicts all block
positions in parallel, the sampler must run many refinement forwards to get
acceptable token quality. In the settings where generations were coherent
enough to verify well, the number of diffusion passes was close to the number
of useful block tokens. The result was a draft-side bottleneck rather than a
hidden draft cost. Our best traditional block-diffusion throughput was only
about $20$ tokens/s, far below the AR and SSD reference points in
Table~\ref{tab:ssd-reference}.

\paragraph{DFlash improves the diffusion primitive but not yet the full hybrid.}
The DFlash path removes the main flaw above: it uses fresh target activations
and proposes the block in one draft call. Our current hybrid, however, is not
yet competitive end to end. On the Qwen3-8B/Qwen3-0.6B/DFlash-b16 prototype,
a 32-token smoke run measured:

\begin{table}[t]
  \caption{Current DFlash+S$^2$D prototype diagnostics. This run used Qwen3-8B
  target, Qwen3-0.6B AR drafter, and \texttt{Qwen3-8B-DFlash-b16}; it was a
  small diagnostic run, not the full 512-prompt evaluation.}
  \label{tab:prototype}
  \centering
  \footnotesize
  \begin{tabular}{@{}lr@{}}
    \toprule
    Metric & Value \\
    \midrule
    Throughput & 11.68 tok/s \\
    Average accepted length & 3.30 tokens \\
    S$^2$D cache hit rate & 0.60 \\
    DFlash seed latency & 34.48 ms \\
    SSD tree-decode latency & 79 ms \\
    \bottomrule
  \end{tabular}
\end{table}

These numbers should not be interpreted as evidence that DFlash is slow.
Rather, they show that our first hybrid pays too much systems overhead for the
small verifier used in the test. The DFlash seed itself is tens of
milliseconds, while the SSD tree-decode path is even larger. With an 8B
verifier on busy A6000s, target verification is not slow enough for the
separate-GPU draft path, tensor transfers, cache bookkeeping, and tree
construction to be hidden.

\paragraph{Small verifiers are the wrong stress test.}
S$^2$D benefits when the target verifier is much slower than the draft path.
If the verifier is small, the latency gap between verifier and drafter shrinks
and fixed costs dominate: inter-GPU tensor transfers, process synchronization,
PyTorch dispatch, CUDA graph warmup, cache-key matching, and tree-cache
population. This explains why the Qwen3-32B SSD speedup reported by
Kumar et al.~\cite{saguaro2024} is lower than the Llama-70B speedup, and why
our Qwen3-8B prototype is an especially unfavorable testbed.

\paragraph{Compute limited the representative test.}
The decisive experiment is DFlash+S$^2$D on a larger verifier, such as
Qwen3-32B, Qwen3-Coder-30B-A3B, Qwen3.6-35B-A3B, or Llama-70B, with the
DFlash and AR draft paths isolated on a separate GPU. We were not able to run
that full experiment within the available GPU credits and shared-server
availability. This is a practical limitation, not just an engineering
inconvenience: the hypothesis specifically concerns latency hiding at large
target-model scale, and small smoke runs cannot validate it.

\subsection{Discussion}

The negative result is still useful. A generic block-diffusion model is not a
drop-in win for S$^2$D because iterative denoising is too expensive. DFlash is
the more plausible route because it collapses block drafting to one
target-conditioned pass and achieves high acceptance in published serving
experiments. Our hybrid design also preserves the strongest part of S$^2$D:
cache hits still return precomputed AR branches immediately. The open question
is whether DFlash can make misses cheap enough on a verifier large enough that
the target pass dominates all other overhead.
